A turn is required at every move: an immediate reversal revisits a square, so horizontal and vertical moves alternate. Colour a cell by the parities of its row and column, using four colours. Along such a path the four colours repeat in order, so every four consecutive visited cells include one of each colour. The colour with both coordinates odd (using indices zero through eight) has only $4\cdot4=16$ cells. Therefore at most $4\cdot16+3=67$ cells can be visited, giving at most 66 moves. The diagram numbers a valid path from zero to 66 with no repeated cell and alternating axes at each move. Thus the maximum is $\boxed{66}$.
\begin{asy}
texpreamble("\usepackage{amsmath,amssymb}");
size(430);
defaultpen(fontsize(9pt));
pair[] p={(1.5,0.5),(0.5,0.5),(0.5,1.5),(1.5,1.5),(1.5,2.5),(0.5,2.5),(0.5,3.5),(1.5,3.5),(1.5,4.5),(0.5,4.5),(0.5,5.5),(1.5,5.5),(1.5,6.5),(0.5,6.5),(0.5,7.5),(1.5,7.5),(1.5,8.5),(2.5,8.5),(2.5,7.5),(3.5,7.5),(3.5,8.5),(4.5,8.5),(4.5,7.5),(5.5,7.5),(5.5,8.5),(6.5,8.5),(6.5,7.5),(7.5,7.5),(7.5,6.5),(6.5,6.5),(6.5,5.5),(7.5,5.5),(7.5,4.5),(8.5,4.5),(8.5,3.5),(7.5,3.5),(7.5,2.5),(8.5,2.5),(8.5,1.5),(7.5,1.5),(7.5,0.5),(6.5,0.5),(6.5,1.5),(5.5,1.5),(5.5,0.5),(4.5,0.5),(4.5,1.5),(3.5,1.5),(3.5,2.5),(2.5,2.5),(2.5,3.5),(3.5,3.5),(3.5,4.5),(2.5,4.5),(2.5,5.5),(3.5,5.5),(3.5,6.5),(4.5,6.5),(4.5,5.5),(5.5,5.5),(5.5,4.5),(4.5,4.5),(4.5,3.5),(5.5,3.5),(5.5,2.5),(6.5,2.5),(6.5,3.5)};for(int i=0;i<=9;++i){draw((i,0)--(i,9),gray);draw((0,i)--(9,i),gray);}for(int i=1;i<p.length;++i)draw(p[i-1]--p[i],blue+linewidth(1));for(int i=0;i<p.length;++i){fill(circle(p[i],.20),white);label(string(i),p[i],fontsize(8pt));}
\end{asy}
